\documentclass[10pt, a4paper]{article}

% ── PACKAGES ──────────────────────────────────────────────────
\usepackage[utf8]{inputenc}
\usepackage[T1]{fontenc}
\usepackage{geometry}
\geometry{
  top=2.0cm,
  bottom=2.0cm,
  left=2.2cm,
  right=2.2cm
}
\usepackage{hyperref}
\hypersetup{
  colorlinks=true,
  urlcolor=blue,
  linkcolor=blue,
  citecolor=blue,
  pdftitle={Pre-Registration v2.0: Plant Architecture Inversion via
            Coherence Gradient Field Specification},
  pdfauthor={Eric Robert Lawson / OrganismCore},
  pdfsubject={Pre-registered experimental protocol and predictions},
  pdfkeywords={plant architecture, coherence gradient, tropism,
               attractor geometry, hydrotropism, phototropism,
               skotomorphogenesis, reticulated foam, cone membrane,
               pre-registration, OrganismCore}
}
\usepackage{microtype}
\usepackage{parskip}
\usepackage{booktabs}
\usepackage{array}
\usepackage{tabularx}
\usepackage{xcolor}
\usepackage{mdframed}
\usepackage{enumitem}
\usepackage{titlesec}
\usepackage{lmodern}
\usepackage{fancyhdr}
\usepackage{lastpage}

% ── COLOURS ───────────────────────────────────────────────────
\definecolor{headergreen}{RGB}{20,100,60}
\definecolor{rulecolor}{RGB}{20,100,60}
\definecolor{claimbox}{RGB}{20,100,60}
\definecolor{notebox}{RGB}{240,248,240}
\definecolor{warningbox}{RGB}{255,248,220}

% ── HEADER / FOOTER ───────────────────────────────────────────
\pagestyle{fancy}
\fancyhf{}
\renewcommand{\headrulewidth}{0.4pt}
\fancyhead[L]{\small\color{headergreen}
  OrganismCore Pre-Registration v2.0}
\fancyhead[R]{\small\color{headergreen}
  Plant Architecture Inversion --- 2026-03-18}
\fancyfoot[C]{\small Page \thepage\ of \pageref{LastPage}}

% ── SECTION FORMAT ────────────────────────────────────────────
\titleformat{\section}
  {\large\bfseries\color{headergreen}}
  {\thesection.}{0.5em}{}[
  \vspace{-0.3em}
  \textcolor{rulecolor}{\rule{\linewidth}{0.6pt}}]
\titleformat{\subsection}
  {\normalsize\bfseries\color{headergreen}}
  {\thesubsection.}{0.5em}{}

\setlength{\parskip}{0.4em}
\setlength{\parindent}{0pt}

% ══════════════════════════════════════════════════════════════
\begin{document}

% ── TITLE BLOCK ───────────────────────────────────────────────
\begin{center}
  {\LARGE\bfseries\color{headergreen}
  PRE-REGISTRATION --- VERSION 2.0}\\[0.4em]
  {\Large\bfseries
  Plant Architecture Inversion via\\
  Coherence Gradient Field Specification}\\[0.5em]
  {\large
  Architectural Inversion of \textit{Raphanus sativus}
  Produced by Pre-Specified Inversion\\
  of the Hydrotropic and Phototropic
  Coherence Gradient Field Before Germination\\
  Five-Scenario Two-Chamber Cone Membrane System ---
  Operative Protocol v2.3}\\[0.6em]
  \noindent\textcolor{rulecolor}{\rule{\linewidth}{1pt}}\\[0.4em]
  {\normalsize
  \textbf{Author:} Eric Robert Lawson\\
  \textbf{Affiliation:} OrganismCore (Independent Research)\\
  \textbf{ORCID:}
  \href{https://orcid.org/0009-0002-0414-6544}
  {0009-0002-0414-6544}\\
  \textbf{Contact:}
  \href{mailto:OrganismCore@proton.me}
  {OrganismCore@proton.me}\\
  \textbf{Repository:}
  \href{https://github.com/Eric-Robert-Lawson/attractor-oncology}
  {\texttt{github.com/Eric-Robert-Lawson/attractor-oncology}}\\
  \textbf{Pre-registration timestamp:} 2026-03-18\\
  \textbf{Linked pre-registration DOI (LECA):}
  \href{https://doi.org/10.5281/zenodo.18986790}
  {\texttt{10.5281/zenodo.18986790}}\\
  \textbf{Supersedes:} v1.0
  (DOI: \href{https://doi.org/10.5281/zenodo.19040399}
  {\texttt{10.5281/zenodo.19040399}},
  timestamp: 2026-03-13)\\
  \textbf{Document version:} 2.0\\
  \textbf{License:}
  Creative Commons Attribution 4.0 International (CC BY 4.0)}\\[0.4em]
  \noindent\textcolor{rulecolor}{\rule{\linewidth}{1pt}}
\end{center}

\vspace{0.3em}

% ── VERSION NOTE ──────────────────────────────────────────────
\begin{mdframed}[
  backgroundcolor=warningbox,
  linecolor=rulecolor,
  linewidth=0.8pt,
  innertopmargin=0.5em,
  innerbottommargin=0.5em,
  innerleftmargin=0.8em,
  innerrightmargin=0.8em
]
\textbf{Amendment note v2.0 (2026-03-18):}
This document supersedes Pre-Registration v1.0
(2026-03-13, DOI: \texttt{10.5281/zenodo.19040399}).
The primary prediction P-INVERSION-1 is unchanged.
The experimental apparatus has been substantially upgraded
from the agar-cup system to a two-chamber opaque PVC pipe
system with a cone membrane signal separator, three
substrate types, and one free UV condition.
The number of experimental scenarios has been expanded
from four to five. Scenario 3 in this document is the
direct equivalent of the EXP chamber in v1.0.
All predictions were updated and re-locked at the
timestamp above before any experimental data exists.
No experimental data has been collected under either
the v1.0 or v2.0 protocol at the time of this upload.
\end{mdframed}

\vspace{0.5em}

% ── STATEMENT BOX ─────────────────────────────────────────────
\begin{mdframed}[
  backgroundcolor=claimbox,
  linecolor=claimbox,
  linewidth=0pt,
  innertopmargin=0.6em,
  innerbottommargin=0.6em,
  innerleftmargin=1.0em,
  innerrightmargin=1.0em
]
\color{white}
\textbf{THE PRE-REGISTERED CLAIM}

The organism is not the output of a genome alone.
It is the output of a genome operating within a coherence
gradient field.
Change the field before germination and the genome executes
a different developmental program.
A different architectural form emerges.
Same genome. Different field. Different organism.

\medskip
This pre-registration specifies: the coherence gradient field
conditions that will produce architectural inversion in
\textit{Raphanus sativus} (radish), the precise architectural
outcome predicted for each of five experimental scenarios,
the scoring method, the control conditions required, the
success criteria, and the reproducibility requirements ---
all locked before the experiment is run.

\medskip
\textbf{Primary prediction (P-INVERSION-1):}
A radish seed placed in Scenario 3 --- hydroponic foam
above, UV-A LED embedded in reticulated foam below, complete
external light deprivation --- will germinate with root
growing upward and shoot growing downward.
Architecture Inversion Score (AIS)\,=\,2 confirmed at
T\,=\,72\,h and maintained at T\,=\,168\,h.
Scenario 1 must record AIS\,=\,0 in the same run.
\end{mdframed}

\vspace{1em}

% ── SECTIONS ──────────────────────────────────────────────────

\section{Background and Theoretical Basis}

\subsection{The Standard Model and Its Incompleteness}

Contemporary biology operates on the model:
\textbf{Genome $\rightarrow$ Organism.}
The genome is treated as the blueprint.
To change the organism, one changes the genome
(CRISPR, directed evolution, synthetic biology).

This model is incomplete.

The complete model is:
\textbf{Genome $+$ Coherence Gradient Field $\rightarrow$ Organism.}

The genome contains \textit{build capacity} ---
the capacity to build every structure the organism has built
across its evolutionary history, and potentially more.
The coherence gradient field is the \textit{instruction set}
that specifies which programs to express, where, and toward
which attractor the developmental program navigates.

The organism that emerges is the resolution of the genome's
build capacity within the field's instruction set.

\subsection{Plant Architecture as Field Navigation}

Standard plant architecture --- root growing downward,
shoot growing upward --- is not genetically fixed.
It is the attractor the genome resolves to within
the standard coherence gradient field:
water and minerals below, light above, gravity present.

Each axis navigates toward its coherence gradient:
\begin{itemize}[itemsep=0.1em]
  \item \textbf{Root:} Navigates toward water and minerals
    (positive hydrotropism, MIZ1/MIZ2 pathway),
    darkness (negative phototropism via phot1/phot2),
    and gravity (positive gravitropism).
  \item \textbf{Shoot:} Navigates toward light
    (positive phototropism via phot1/phot2),
    CO\textsubscript{2}, and warmth
    (negative gravitropism).
\end{itemize}

Standard architecture is the attractor produced when all
gradients are coherent: root signals (water, minerals,
darkness) point downward; shoot signals (light) point upward.
Gravity confirms both directions.

\subsection{The Signal Deprivation Principle}

The central engineering condition of this experiment
(Scenarios 2--5) is \textbf{complete external light
deprivation.} The UV-A LED embedded within the foam
substrate is the \textit{only} photonic signal in the
coherence field of each chamber.

When a germinating seedling develops in complete darkness
it enters \textbf{skotomorphogenesis} --- the developmental
program for light-seeking under total light deprivation.
In this state:
\begin{itemize}[itemsep=0.1em]
  \item Hypocotyl elongation is maximal.
  \item All photoreceptors are primed at maximum sensitivity.
  \item There is zero ambient photonic noise.
\end{itemize}

Under these conditions the absolute efficiency of UV-A as
a phototropic driver is irrelevant. The embedded UV-A
source is not competing against anything.
The signal-to-noise ratio is effectively infinite.
UV-A activates phototropins (phot1, phot2).
In a completely dark chamber with a single directional
UV-A source, a skotomorphogenic seedling responds to the
only photonic input in its entire coherence field.

Scenario 1 is the explicit exception: UV-A is freely
available in the S1 air space with zero penetration cost.
This is the true normal control --- normal conditions,
normal architecture confirmed.

\subsection{The Geometric Prediction}

If plant architecture is a field-navigated attractor
rather than a genetically fixed output, then:

\textbf{Inverting the coherence gradient field before
germination should produce an inverted architectural attractor.}

Placing hydroponic substrate above the seed and UV-A
embedded substrate below the seed --- before germination
begins --- should cause the root to navigate upward
(toward water, suppressing gravitropism via MIZ1-mediated
amyloplast degradation) and the shoot to navigate downward
(toward the only photonic signal in the field).

This prediction is geometrically derived from the attractor
framework \textit{before} the experiment is run.
The conditions are specified from the geometry.
The outcome is predicted from the conditions.
The experiment confirms or falsifies the geometry.

\subsection{Empirical Basis for Individual Mechanisms}

All component mechanisms are independently established
in the published literature.
No mechanism claimed here is novel in isolation.
The novel claim is their combination in a pre-specified
five-scenario structured experiment producing simultaneous
full architectural inversion from a locked prediction.

\begin{itemize}[itemsep=0.2em]
  \item \textbf{Hydrotropism overriding gravitropism, radish:}
    Takahashi et al.\ (2003), \textit{Plant Physiology} 132:
    805--810.
    Radish (\textit{Raphanus sativus}) roots grow upward toward
    water. Mechanism: MIZ1-mediated amyloplast degradation
    suppresses gravity sensing. Confirmed in both radish
    and \textit{Arabidopsis}.

  \item \textbf{MIZ1 suppresses gravitropism under drought:}
    Zhang et al.\ (2025), \textit{PNAS} 122(20): e2427315122.
    Under drought, roots ignore gravity and grow toward water
    (including upward). MIZ1 suppresses PIN auxin transporter
    activity that drives gravitropism.
    Confirmed in laboratory, clinostat, and field conditions.
    Direct relevance: in Scenario 3, the root faces a
    drought-like condition --- all water is above, none below.
    MIZ1-mediated suppression of gravity sensing is
    the predicted operative mechanism.

  \item \textbf{Root negative phototropism to UV-A/blue:}
    Silva-Navas et al.\ (2016), \textit{PNAS} 113(11):
    3119--3124.
    Roots actively avoid blue and UV-A wavelengths.
    Mechanism: asymmetric auxin redistribution via phot1/phot2.
    In Scenario 3, this drives root away from UV-A below,
    adding a second mechanism pointing root upward.

  \item \textbf{Hydrotropism overrides root negative
    phototropism:}
    When water and light are in competition, the hydrotropic
    signal dominates when the water gradient is sufficient
    (Kaneyasu et al., 2007).
    In Scenario 3, root negative phototropism (away from
    UV-A below) and positive hydrotropism (toward water
    above) both point upward. No conflict. Both mechanisms
    are cooperative.

  \item \textbf{Shoot positive phototropism toward buried light:}
    Christie (2007), \textit{Annual Review of Plant Biology}
    58: 21--45.
    Shoot phototropism is always positive regardless of
    light direction. When light is below, the hypocotyl
    bends downward. Mediated by phot1/phot2.
    In Scenario 3, the only photonic signal is UV-A
    embedded in Type B foam below. The shoot bends toward it.

  \item \textbf{Root positive phototropism unmasked in
    microgravity:}
    Kiss et al.\ (2012), \textit{American Journal of Botany}
    and NASA TROPI ISS experiments.
    In microgravity, Arabidopsis roots show positive
    phototropism toward blue light --- a capacity completely
    masked by gravitropism on Earth.
    Direct framework relevance: the genome contains programs
    that are suppressed by the standard coherence gradient
    field and become expressed when a component of that
    field is removed or overridden.

  \item \textbf{Photosynthesis from geothermal light
    (buried source):}
    Beatty et al.\ (2005), \textit{PNAS} 102(26):
    9306--9310.
    Green sulfur bacteria at 2,500\,m depth perform
    photosynthesis using geothermal light from below.
    Confirmed analogue: photonic navigation toward
    a buried non-solar source is a real biological
    phenomenon across kingdoms.

  \item \textbf{Radiotropism toward embedded radiation:}
    Dadachova et al.\ (2007), \textit{PLoS ONE}.
    Melanized fungi at Chernobyl grow toward ionizing
    radiation and use it as an energy source.
    Confirmed analogue: organisms navigate toward and
    harvest embedded non-standard energy sources
    under field deprivation of standard sources.
\end{itemize}

\vspace{0.5em}

\section{Pre-Registered Predictions}

All predictions below are locked at the timestamp
of this document's upload to Zenodo.
No post-hoc modifications will be made to predictions.
Amendments, if any, will be registered as separate
versioned documents with new DOIs.

\subsection{Architecture Inversion Score Definition}

The Architecture Inversion Score (AIS) is the primary
quantitative outcome measure, defined as follows.
Angles are measured from the vertical
(straight down = 0\textdegree, straight up = 180\textdegree)
using a protractor held against the transparent section
of the chamber wall at each observation timepoint.

\medskip
\begin{tabular}{p{1.2cm}p{4.8cm}p{5.5cm}}
\toprule
\textbf{AIS} & \textbf{Morphological Criterion} &
\textbf{Interpretation} \\
\midrule
2 &
Root angle $>135$\textdegree{} AND shoot angle
$<45$\textdegree{} from vertical downward
(both axes fully inverted). &
Full architectural inversion.
P-INVERSION-1 confirmed. \\[0.4em]
1 &
One axis inverted ($>90$\textdegree{} from expected),
other axis partially or not inverted. &
Partial inversion.
One gradient dominant, other insufficient. \\[0.4em]
0 &
Root angle $<45$\textdegree{} (downward) AND
shoot angle $>135$\textdegree{} (upward). &
Standard architecture.
No inversion. \\[0.4em]
$-1$ &
Both axes in same unexpected direction
(both up or both down). &
Disorientation.
Mechanism failure or contamination. \\
\bottomrule
\end{tabular}

\subsection{Primary Prediction: P-INVERSION-1}

\begin{mdframed}[
  backgroundcolor=notebox,
  linecolor=rulecolor,
  linewidth=0.8pt,
  innertopmargin=0.5em,
  innerbottommargin=0.5em,
  innerleftmargin=0.8em,
  innerrightmargin=0.8em
]
\textbf{P-INVERSION-1 (Primary):}
In \textbf{Scenario 3} (S1: Type A hydroponic foam above;
S2: Type B UV-only foam below; cone apex down; complete
external light deprivation), \textit{Raphanus sativus}
will germinate and grow with:
\begin{itemize}[itemsep=0.1em, topsep=0.2em]
  \item Root directed \textbf{upward}
    (angle $>135$\textdegree{} from vertical downward)
    at T\,=\,72\,h.
  \item Measurable root penetration depth into S1
    Type A hydroponic foam recorded at T\,=\,72\,h.
  \item Shoot directed \textbf{downward}
    (angle $<45$\textdegree{} from vertical downward)
    at T\,=\,72\,h.
  \item Measurable shoot penetration depth into S2
    Type B UV foam recorded at T\,=\,72\,h.
  \item Architecture Inversion Score
    \textbf{AIS\,=\,2} at T\,=\,72\,h,
    maintained at T\,=\,120\,h and T\,=\,168\,h.
\end{itemize}
\end{mdframed}

\subsection{Scenario Predictions: P-INVERSION-2
through P-INVERSION-5}

\begin{mdframed}[
  backgroundcolor=notebox,
  linecolor=rulecolor,
  linewidth=0.8pt,
  innertopmargin=0.5em,
  innerbottommargin=0.5em,
  innerleftmargin=0.8em,
  innerrightmargin=0.8em
]
\textbf{P-INVERSION-2 (Scenario 1 --- True Normal Control):}
S1: UV-A LED free in air space (zero penetration cost).
S2: Type A hydroponic foam. Cone apex up.
\textit{Expected:} AIS\,=\,0 at all timepoints.
Root grows downward into S2 Type A foam.
Shoot grows upward into free UV-lit air without
penetrating any foam barrier.
Standard architecture confirmed.
This scenario is the validity gate for every run ---
if AIS\,$\ne$\,0 in Scenario 1, the run is invalid
regardless of any other result.

\medskip
\textbf{P-INVERSION-3 (Scenario 2 --- UV separated above,
hydroponic below):}
S1: Type B UV-only foam (UV-A embedded, no hydroponic).
S2: Type A hydroponic foam. Cone apex up.
\textit{Expected:} Root grows downward into S2 Type A.
Shoot grows upward and penetrates S1 Type B foam to
reach embedded UV-A. Normal architecture maintained
but active shoot penetration upward through UV foam required.
AIS\,=\,0 at T\,=\,72\,h with measurable shoot penetration
into S1 Type B.
This scenario isolates the phototropic penetration
component without hydrotropic competition.
Provides the inversion pair baseline against Scenario 3.

\medskip
\textbf{P-INVERSION-4 (Scenario 4 --- All signals above,
void below):}
S1: Type C foam (hydroponic + UV-A co-located). Void in S2.
Cone apex up.
\textit{Expected:} Shoot grows upward penetrating S1 Type C
to reach all co-located signals simultaneously.
Root grows downward into void or is indeterminate.
AIS\,=\,0 or indeterminate.
Maximum coherent field above, no competing signal below.

\medskip
\textbf{P-INVERSION-5 (Scenario 5 --- All signals below,
void above):}
S2: Type C foam (hydroponic + UV-A co-located). Void in S1.
Cone apex down.
\textit{Expected:} Full inversion. Shoot grows downward
penetrating S2 Type C to reach all co-located signals.
Root grows upward into void or is indeterminate.
AIS\,=\,2 at T\,=\,72\,h.
Maximum coherent field below, no competing signal above.
This is the strongest possible inversion signal
(no competing attractor in opposite direction).
Comparison of Scenario 5 (AIS\,=\,2, no competition)
vs Scenario 3 (AIS\,=\,2, competing hydroponic above)
tests the effect of competing signal on inversion
stability.
\end{mdframed}

\vspace{0.5em}

\section{Experimental Design --- v2.3 Apparatus}

\subsection{Overview}

Five scenarios are run simultaneously in the same
environment from the same seed batch and same
substrate preparation batch.
The only variable between scenarios is the configuration
of the coherence gradient field.
All 5 scenarios use the same two-chamber cone membrane
apparatus with the same UV-A LED product.

\subsection{The Two-Chamber Cone Membrane System}

Each experimental unit consists of two opaque PVC pipe
sections (50\,mm internal diameter, 150\,mm length each)
joined at a cone membrane signal separator.

\textbf{The cone membrane} is a plastic funnel wrapped in
metallized Mylar film (aluminized emergency blanket),
serving two barrier functions:
\begin{enumerate}[itemsep=0.1em]
  \item \textbf{Liquid barrier:} prevents hydroponic
    solution migrating between S1 (top) and S2 (bottom).
  \item \textbf{UV and light barrier:} prevents UV-A
    from the substrate in one chamber crossing into the
    other, and prevents external light entering at the
    joint.
\end{enumerate}

A hydroponic net cup seats at the cone apex.
The germinated seed is placed in the net cup with
radicle extending through the net cup mesh.
The cone orientation (apex up or apex down) determines
which chamber the radicle faces and which chamber the
emerging hypocotyl faces.

\textbf{Cone orientation rule:}
The apex always points toward the void or least active
substrate. The wide opening always faces the most
active substrate.

\subsection{Substrate Types}

Three substrate types and one free UV condition are
used across the five scenarios:

\begin{itemize}[itemsep=0.3em]
  \item \textbf{Free UV condition (Scenario 1 S1 only):}
    UV-A LED strip mounted free on inner face of S1 end cap,
    facing downward into S1 air space.
    No foam substrate in S1. UV-A freely available throughout
    S1 air space. Zero penetration cost.
    No water or minerals.

  \item \textbf{Type A --- Hydroponic foam, no UV:}
    Reticulated polyurethane foam 20\,ppi, saturated with
    hydroponic mineral solution (1/4 to 1/2 strength).
    No UV-A LED present at any depth.
    Signals: water and minerals only.

  \item \textbf{Type B --- UV-only foam, no hydroponic:}
    Reticulated polyurethane foam 20\,ppi, dry or minimally
    moist (plain distilled water only, no nutrients), with
    UV-A LED strip embedded 10--20\,mm depth from the
    seed-facing surface.
    Signals: UV-A radiation only.
    Seedling must physically penetrate foam to reach UV-A.

  \item \textbf{Type C --- Combined foam, hydroponic + UV:}
    Reticulated polyurethane foam 20\,ppi, saturated with
    hydroponic mineral solution, with UV-A LED strip embedded
    10--20\,mm depth from the seed-facing surface.
    Signals: water, minerals, and UV-A co-located.
    Seedling must penetrate foam to reach all signals.
\end{itemize}

\subsection{The Five Scenarios}

\begin{tabularx}{\linewidth}{p{1.1cm}p{3.0cm}p{3.0cm}p{1.1cm}X}
\toprule
\textbf{ID} & \textbf{S1 (top)} & \textbf{S2 (bottom)} &
\textbf{Cone} & \textbf{Purpose and prediction} \\
\midrule
S1 &
Free UV in air &
Type A hydroponic &
Apex up &
True normal control.
AIS\,=\,0 required.
Validity gate. \\[0.3em]
S2 &
Type B (UV only) &
Type A hydroponic &
Apex up &
UV separated above.
Normal architecture + shoot penetrates Type B upward.
AIS\,=\,0.
Inversion pair baseline. \\[0.3em]
\textbf{S3} &
\textbf{Type A hydroponic} &
\textbf{Type B (UV only)} &
\textbf{Apex dn} &
\textbf{PRIMARY CLAIM.}
Full inverted field.
AIS\,=\,2 predicted.
Root up into Type A.
Shoot down penetrating Type B. \\[0.3em]
S4 &
Type C (all signals) &
Empty void &
Apex up &
All signals above. Normal + maximum coherent field.
AIS\,=\,0. \\[0.3em]
S5 &
Empty void &
Type C (all signals) &
Apex dn &
All signals below. Maximum inversion signal.
AIS\,=\,2 predicted.
No competing attractor above. \\
\bottomrule
\end{tabularx}

\subsection{Materials}

\begin{tabular}{p{4.5cm}p{2.2cm}p{4.5cm}}
\toprule
\textbf{Item} & \textbf{Cost (USD)} & \textbf{Source} \\
\midrule
\textit{Raphanus sativus} seeds
(Cherry Belle) &
\$2--4 & Garden store, Amazon \\
Reticulated polyurethane foam 20\,ppi &
\$5--10 & Aquarium supply, Amazon \\
Hydroponic mineral solution
(seedling dilution, 1/4--1/2 strength) &
\$5--12 & Garden centre, hydro supplier \\
Opaque PVC pipe 50\,mm ID, 150\,mm lengths $\times$10 &
\$10--20 & Plumbing / hardware store \\
PVC end caps $\times$10 &
\$5--10 & Plumbing / hardware store \\
UV-A LED strip $\sim$365\,nm &
\$8--15 & Amazon \\
Metallized Mylar film (emergency blanket) &
\$1--3 & Camping supply, dollar store \\
Plastic funnels (50\,mm rim) $\times$5 &
\$3--6 & Kitchen supply \\
Hydroponic net cups 16--25\,mm $\times$5 &
\$1--3 & Hydro supplier, Amazon \\
Aquarium silicone sealant &
\$5--8 & Aquarium store, hardware \\
Distilled water &
\$1--2 & Grocery store \\
Red-light torch ($>$650\,nm) &
\$5--15 & Astronomy / photography \\
\midrule
\textbf{Total} & \textbf{\$51--108} & \\
\bottomrule
\end{tabular}

\subsection{Hydroponic Mineral Solution Composition}

Option A (from salts):
\begin{itemize}[itemsep=0.1em]
  \item 1\,mM KNO\textsubscript{3}
  \item 0.5\,mM KH\textsubscript{2}PO\textsubscript{4}
  \item 0.2\,mM MgSO\textsubscript{4}
  \item pH 5.8--6.5, adjusted with dilute KOH if needed
  \item Prepared in distilled water
\end{itemize}

Option B (simplified):
Commercially available balanced hydroponic nutrient
solution at 1/4 to 1/2 manufacturer's recommended
seedling dilution, prepared in distilled water.
Adequate ionic content for hydrotropism signalling.

\subsection{Seed Germination Protocol}

Seeds are germinated on damp paper towel in a sealed
ziplock bag in complete darkness at 20--24\textdegree C
for 48\,hours before placement.
Target state at placement: radicle 3--8\,mm, intact,
no kinking.
All germination and placement steps under red-light only
($>$650\,nm) to preserve skotomorphogenic state.
The seedling must not be exposed to white light,
blue light, or UV at any point from germination
through placement to T\,=\,0.

\subsection{Light Exclusion Protocol}

Complete external light deprivation for Scenarios 2--5
is the load-bearing engineering constraint.

\begin{enumerate}[itemsep=0.2em]
  \item Use opaque black PVC pipe (preferred) OR wrap
    full exterior of each chamber unit in black opaque
    tape or black plastic wrap with no gaps.
  \item All ventilation ports fitted with reticulated
    foam plug light traps (20--25\,mm depth, 20\,ppi).
    Test each port: shine torch at outer face of foam
    plug; confirm no light visible from interior side.
  \item Cone membrane joint sealed with continuous
    aquarium silicone bead. Light-seal test after curing:
    place torch inside assembled chamber; check for light
    leaks at joint from outside.
  \item At each observation timepoint: red-light only
    for all interior inspection. Maximum 2 minutes
    light exposure per chamber per timepoint.
    Re-wrap or re-enclose immediately after photography.
    Record duration of light exposure at each timepoint.
\end{enumerate}

Scenario 1 S1 is exempt from external light exclusion
--- the free UV LED is the design condition of that
chamber. Scenario 1 S2 should be opaque.

\subsection{Observation Protocol}

\begin{tabular}{p{1.8cm}p{10.0cm}}
\toprule
\textbf{Timepoint} & \textbf{Observation} \\
\midrule
T\,=\,24\,h &
Germination check. Record radicle length and direction
per scenario. Photograph from side under red-light.
Confirm UV-A LED active in all powered scenarios.
Re-seal and re-wrap Scenarios 2--5 immediately. \\[0.3em]
T\,=\,48\,h &
Directional growth assessment.
Record root and shoot direction per scenario
(up/down/horizontal/indeterminate).
Measure penetration depth into substrate (mm).
Preliminary AIS. Re-seal and re-wrap. \\[0.3em]
T\,=\,72\,h &
\textbf{Primary outcome point.}
Measure root and shoot angles with protractor.
Measure substrate penetration depth (mm).
Record whether root or shoot tip has reached
LED embedding depth in Type B or Type C foam.
Calculate AIS for all scenarios.
Photograph from side and above with scale ruler.
Re-seal and re-wrap. \\[0.3em]
T\,=\,120\,h &
Secondary confirmation.
AIS re-measured. Stability of inversion assessed.
Note any gravitropic correction (inversion at T\,=\,72
but recovered at T\,=\,120). \\[0.3em]
T\,=\,168\,h &
Final outcome. Photograph all scenarios together on
white background. Side-by-side comparison of Scenario 3
vs Scenario 1. Record final AIS for all scenarios. \\
\bottomrule
\end{tabular}

\vspace{0.5em}

\section{Success Criteria and Falsification}

\subsection{Confirmation of P-INVERSION-1}

P-INVERSION-1 is confirmed if and only if all of the
following are met in the same experimental run:
\begin{enumerate}[itemsep=0.2em]
  \item Scenario 3: AIS\,=\,2 at T\,=\,72\,h
    \textbf{and} maintained at T\,=\,168\,h.
  \item Scenario 3: measurable shoot penetration into
    S2 Type B UV foam recorded at T\,=\,72\,h.
    Directional growth alone without foam penetration
    does not confirm active navigation toward the
    UV-A signal.
  \item Scenario 1: AIS\,=\,0 at all timepoints.
    If Scenario 1 records AIS\,$\ne$\,0, the run is
    invalid regardless of Scenario 3 result.
  \item Complete light exclusion confirmed and logged
    for Scenarios 2--5 at every timepoint.
    Runs with unlogged or compromised light exclusion
    must be reported as flagged and analysed separately
    from primary results.
\end{enumerate}

\textbf{Reproducibility threshold:}
P-INVERSION-1 requires confirmation across
\textbf{minimum 3 independent seeds} in
\textbf{minimum 2 independent experimental runs}
from separate seed batches before the result is
considered reportable as confirmed.

\subsection{Qualified Results}

\textbf{Gravitropic correction (AIS\,=\,2 at T\,=\,72,
AIS\,=\,0 at T\,=\,120):}
This is a qualified result, not a failure.
It confirms the coherence field inversion succeeded in
redirecting architecture during the developmental window
and that gravitropic correction occurs after window
closure. Report full angle data at all timepoints.
This is data about the competition between the coherence
gradient field and gravitropism over developmental time.

\textbf{AIS\,=\,1 in Scenario 3:}
Partial inversion. One axis inverted, one insufficient.
Not a confirmation of P-INVERSION-1.
Report as qualified result. Investigate whether
signal strength (LED embedding depth, intensity)
or foam moisture level requires adjustment.

\subsection{Falsification Criteria}

The geometric prediction is falsified if, across
minimum 3 seeds in 2 independent runs with confirmed
light exclusion and confirmed Scenario 1 AIS\,=\,0:
\begin{itemize}[itemsep=0.1em]
  \item Scenario 3 AIS\,=\,0 at T\,=\,72\,h in all runs
    (gravity dominates over both combined coherence
    gradients at this field strength).
  \item Scenario 3 AIS\,=\,$-1$ in all runs
    (developmental confusion rather than inversion).
\end{itemize}

A single run failure does not falsify the prediction.
Falsification requires consistent failure across the
minimum reproducibility threshold with all controls
confirmed valid.

\vspace{0.5em}

\section{Reproducibility Requirements}

For any replication of this experiment to be considered
valid and comparable to this pre-registration, the
following minimum documentation must be provided:

\begin{enumerate}[itemsep=0.2em]
  \item Species, cultivar, and seed source (supplier,
    batch number, viability test result).
  \item Substrate specification: foam ppi rating, supplier;
    nutrient solution product/formulation and dilution
    ratio; LED product name, stated wavelength (nm),
    embedding depth (mm from seed-facing surface),
    moisture condition of Type B foam.
  \item Chamber specification: pipe material, inner
    diameter, section length; cone membrane material
    and construction; net cup dimensions.
  \item Light exclusion method and test result for each
    chamber at T\,=\,0 (confirmed light-tight Y/N per
    chamber).
  \item Red-light torch wavelength and confirmation that
    white/blue/UV light was not used during any
    observation.
  \item Temperature throughout experiment.
  \item All 5 scenarios run simultaneously from T\,=\,0
    in the same environment.
  \item AIS scores and substrate penetration depths (mm)
    at T\,=\,72\,h, T\,=\,120\,h, T\,=\,168\,h
    for all 5 scenarios.
  \item UV-A LED active status confirmed at each
    timepoint for all powered scenarios.
  \item Duration of light exposure per chamber per
    timepoint recorded.
  \item Minimum one photograph per scenario per
    timepoint, labelled with scenario ID, timepoint,
    and date/time.
  \item Any interventions or failure events recorded
    with timestamps and data-compromise flags.
\end{enumerate}

\vspace{0.5em}

\section{The Comparison Structure}

The five-scenario design is not five independent
experiments. It is a structured reasoning space over
two variables --- UV direction and signal separation ---
that makes the primary result (Scenario 3) unambiguous
by eliminating alternative explanations.

\begin{itemize}[itemsep=0.2em]
  \item \textbf{S1 vs S2:} UV free in air (zero
    penetration cost) vs UV in Type B foam above
    (high penetration cost). Same hydroponic below.
    Tests penetration cost alone.
  \item \textbf{S2 vs S3:} Primary inversion pair.
    Structurally identical --- same substrate types,
    same signal strengths, same penetration requirements.
    Only direction flips. Cone flips with it.
    If architecture flips, the field is the author.
  \item \textbf{S3 vs S5:} Inversion with competing
    hydroponic above vs inversion with nothing above.
    Tests signal competition in the inversion result.
  \item \textbf{S4 vs S5:} All signals above vs all
    signals below. Maximum coherent field.
    Tests whether full inversion occurs when the signal
    is unambiguous and uncontested.
\end{itemize}

\vspace{0.5em}

\section{Relationship to the Broader Framework}

This experiment is one component of a larger
pre-registered research program (OrganismCore,
DOI: \href{https://doi.org/10.5281/zenodo.18986790}
{\texttt{10.5281/zenodo.18986790}}) testing
the geometric claim that organisms are the output
of genomes operating within coherence gradient fields,
and that changing the field changes the organism.

The plant inversion experiment is the first
\textit{prospective} test of this claim ---
a pre-specified prediction confirmed by experiment
rather than a retrospective explanation of
already-observed phenomena.

The retrospective confirmation record
(confirmed cases across multiple kingdoms including
hydrotropism overriding gravitropism in radish,
root positive phototropism in microgravity,
photosynthesis from geothermal buried light,
and radiotropism toward embedded radiation sources,
documented in the OrganismCore repository) is available at:
\href{https://github.com/Eric-Robert-Lawson/attractor-oncology}
{\texttt{github.com/Eric-Robert-Lawson/attractor-oncology}}

The oncological application of the coherence gradient
field framework --- that cancer represents a breakdown
in coherence field maintenance within a tissue, with
cells navigating to the wrong attractor because the
local field has been disrupted --- depends on establishing
the principle that the field authors the form at the
organismal level.
The plant inversion experiment is the foundational
proof of that principle.

\vspace{0.5em}

\section{Declaration}

This document is submitted as a pre-registration.
The experiment described has \textbf{not} been run
at the time of this document's creation and upload.
All predictions are made before any experimental
data exists.
Results, when obtained, will be reported in a
separate document linked to this pre-registration
by DOI.
No modifications will be made to this document's
predictions after upload.
Amendments will be versioned and separately registered
with new DOIs.

\vspace{0.3em}
\textbf{Timestamp:} 2026-03-18\\
\textbf{Document version:} 2.0\\
\textbf{Supersedes:} v1.0 (2026-03-13,
DOI: \texttt{10.5281/zenodo.19040399})\\
\textbf{Author:} Eric Robert Lawson / OrganismCore\\
\textbf{ORCID:}
\href{https://orcid.org/0009-0002-0414-6544}
{0009-0002-0414-6544}

\vspace{1em}
\noindent\textcolor{rulecolor}{\rule{\linewidth}{0.6pt}}
\begin{center}
\footnotesize
\textit{The field came first.\ \
The genome was always listening.\ \
The architecture was never fixed.\ \
The inversion was always possible.}\\[0.3em]
\href{https://github.com/Eric-Robert-Lawson/attractor-oncology}
{\texttt{github.com/Eric-Robert-Lawson/attractor-oncology}}
\ $\cdot$\
\href{https://doi.org/10.5281/zenodo.18986790}
{\texttt{doi.org/10.5281/zenodo.18986790}}
\ $\cdot$\
\href{https://orcid.org/0009-0002-0414-6544}
{\texttt{ORCID: 0009-0002-0414-6544}}
\end{center}

\end{document}
